\SourcePage{296}
\SGAsection{6}{以乘法群的 Néron 模型为核的扩张与双扩张}
\textbf{6.1.} 设 $S$ 为\Term{scheme}{概形}，$T$ 为其\Term{closed-subscheme}{闭子概形}，$Z$ 为\Term{set}{集合}，从而得到 $S$ 上的\Term{truncated-constant-scheme}{截断常值概形} $Z_{S,T}$（5.7）。有一个关于 $Z$ \Term{functoriality}{函子性}的、$S$ 上\Term{etale}{平展}\Term{scheme}{概形}的\Term{canonical}{典范}\Term{homomorphism}{同态}
\[
Z_S\longrightarrow Z_{S,T},\tag{6.1.1}
\]
由（5.7.1）立即可知，它是\Term{surjective-morphism}{满射态射}。由其关于 $Z$ 的\Term{functoriality}{函子性}可知，当 $Z$ 是\Term{group}{群}时，（6.1.1）是 $S$ 上\Term{etale}{平展}\Term{group-scheme}{群概形}的\Term{homomorphism}{同态}。因此其\Term{kernel}{核}也是 $S$ 上的\Term{etale}{平展}\Term{group-scheme}{群概形}，记为 $Z_{S,U}$，并且有关于 $Z$ \Term{functoriality}{函子性}的\Term{canonical}{典范}\Term{exact-sequence}{正合列}
\[
0\longrightarrow Z_{S,U}\longrightarrow Z_S\longrightarrow Z_{S,T}\longrightarrow0.\tag{6.1.2}
\]
由（5.7.1）立即看出，$Z_{S,U}$ 是一族 $S$-\Term{scheme}{概形} $S_n$（$n\in Z$）的和，其中 $n=0$ 时 $S_n=S$，$n\ne0$ 时 $S_n=U$。由 $Z_{S,U}$ 的这个描述容易得到，对 $S$ 上任意\Term{group}{群} $F$，\Term{restriction}{限制}到 $U$ 的映射给出\Term{bijection}{双射}
\SourcePage{297}
\[
\operatorname{Hom}_{\mathrm{gr}}(Z_{S,U},F)\xrightarrow{\sim}
\operatorname{Hom}_{\mathrm{gr}}(Z_U,F|U).\tag{6.1.3}
\]
\textbf{6.1.4.} 以下假定 $Z$ 交换，并研究 $Z_{S,T}$ 的以 $S$ 上给定\Term{commutative-group}{交换群} $G$ 为\Term{kernel}{核}的\Term{commutative-group-extension}{交换群扩张}。由\Term{exact-sequence}{正合列}（6.1.2），这些\Term{group-extension}{扩张}的\Term{category}{范畴}\Term{equivalence}{等价}于如下\Term{category}{范畴}：其\Term{object}{对象}是 $Z_S$ 的以 $G$ 为\Term{kernel}{核}的\Term{group-extension}{扩张} $F$，并附有该\Term{group-extension}{扩张}沿 $Z_{S,U}\to Z_S$ 的\Term{inverse-image}{逆像}的一个\Term{trivialization}{平凡化}。根据（6.1.3），给出这样的\Term{trivialization}{平凡化}，\Term{equivalence}{等价}于给出 $Z_U$ 的以 $G|U$ 为\Term{kernel}{核}的\Term{group-extension}{扩张} $F|U$ 的一个\Term{trivialization}{平凡化}。

以下取 $Z$ 为整数\Term{group}{群} $\mathbf Z$。于是 $\mathbf Z_S$ 是单\Term{generator}{生成元}的\Term{free-group}{自由群}，第VII讲1.4给出 $\mathbf Z_S$ 的以 $G$ 为\Term{kernel}{核}的\Term{group-extension}{扩张}的\Term{category}{范畴}与 $G$-\Term{torsor}{挠子}的\Term{category}{范畴}之间的\Term{equivalence-of-categories}{范畴等价}。对 $\mathbf Z_U$ 和 $G|U$ 应用同一结果，便得到：
\begin{SGAresult}{\Term{lemma}{引理}}{6.2}
设 $S$ 为\Term{scheme}{概形}，$T$ 为其\Term{closed-subscheme}{闭子概形}，$U=S-T$，$G$ 为 $S$ 上的\Term{commutative-group}{交换群}。$S$ 上的\Term{commutative-group}{交换群}中，$\mathbf Z_{S,T}$（定义见5.7）的以 $G$ 为\Term{kernel}{核}的\Term{group-extension}{扩张}所成的\Term{category}{范畴}，\Term{equivalence}{等价}于二元组 $(F,s)$ 的\Term{category}{范畴}，其中 $F$ 是 $G$-\Term{torsor}{挠子}，$s$ 是 $F|U$ 的一个\Term{section}{截面}。若 $\alpha:S\to\mathbf Z_{S,T}$ 是 $\mathbf Z_S$ 在 $S$ 上的\Term{section}{截面} $1$ 的像，则 $\mathbf Z_{S,T}$ 的以 $G$ 为\Term{kernel}{核}的\Term{group-extension}{扩张} $\overline G$ 所对应的\Term{torsor}{挠子}为 $F=\alpha^*(\overline G)$；它所附的 $F|U$ 的\Term{section}{截面}由 $\overline G$ 的\Term{unit-section}{单位截面}在 $U$ 上的\Term{restriction}{限制}给出。
\end{SGAresult}
\SourcePage{298}
\textbf{6.2.1.} 若 $\overline G$ 对应于二元组 $(F,s)$，以 $\mathbf Z_S$ 的值为 $n\in\mathbf Z$ 的\Term{constant}{常值}\Term{section}{截面}定义 $\mathbf Z_{S,T}$ 的\Term{section}{截面} $\alpha_n$，则 $\overline G$ 沿 $\alpha_n$ 的\Term{inverse-image}{逆像} $\alpha_n^*(\overline G)$，作为 $G$-\Term{torsor}{挠子}，\Term{canonical}{典范}地等同于 $G$-\Term{torsor}{挠子} $F$ 的 $n$ 次幂 $F^{(n)}$；由（5.7.1）的第一个关系定义的 $G_U$-\Term{torsor}{挠子}的\Term{canonical-isomorphism}{典范同构}
\[
\alpha_n^*(\overline G)|U\simeq G_U
\]
由 $F_U^{(n)}$ 的\Term{section}{截面} $s^{(n)}$ 确定，后者是给定的 $F_U$ 的\Term{section}{截面} $s$ 的 $n$ 次幂。由这一具体解释，得到 $\overline G$ 在 $S$ 上的\Term{section}{截面}的如下描述：一个这样的\Term{section}{截面} $\overline g$ 由下列数据确定：

 a）$\mathbf Z_{S,T}$ 在 $S$ 上的一个\Term{section}{截面} $\sigma$，它\Term{solvable}{可解}释为 $\mathbf Z_T$ 的一个\Term{section}{截面}，即 $T$ 上取值于 $\mathbf Z$ 的\Term{locally-constant-function}{局部常值函数}；也可把它解释为 $T$ 的一个\Term{partition}{分割}
\[
T=\coprod_{i\in\mathbf Z}T_i,
\]
其中各 $T_i$ 是由 $i\in\mathbf Z$ 编号的两两不交\Term{open-subset}{开集}；

 b）一族\Term{section}{截面}
\[
\overline g_i\in\Gamma(S_i',F^{(i)}|S_i'),\qquad i\in\mathbf Z,
\qquad S_i'=S-T_i',\quad T_i'=\coprod_{\substack{j\in\mathbf Z\\j\ne i}}T_j,
\]
各 $i\in\mathbf Z$ 的\Term{section}{截面} $\overline g_i$ 满足如下相容条件：对 $i,j\in\mathbf Z$，借助这些\Term{torsor}{挠子}的\Term{section}{截面} $s^{(i)}$、$s^{(j)}$，把 $F^{(i)}|U$、$F^{(j)}|U$ 与 $G_U$ 等同后，$\overline g_i$ 与 $\overline g_j$ 在 $S_i'\cap S_j'=U$ 上相同。换言之，必须存在一个显然唯一确定的\Term{section}{截面} $h_{ij}\in\Gamma(U,G|U)$，使得
\[
\overline g_i=h_{ij}s^{(i)},\qquad \overline g_j=h_{ij}s^{(j)}.
\]
由于3.2中对 $\mathbf Z_{S,T}$ 的以 $G$ 为\Term{kernel}{核}的\Term{group-extension}{扩张} $\overline G$ 的描述显然与任意\Term{base-change}{基变换} $S'\to S$ 相容，上述讨论便完整描述了\Term{sheaf}{层} $\overline G$ 作为变元 $S'\in\operatorname{Ob}(\mathrm{Sch}/S)$ 的\Term{functor}{函子}的情形。
\SourcePage{299}
\textbf{6.2.2.} 当 $T$ 是\Term{nonempty}{非空}\Term{connected}{连通}\Term{scheme}{概形} $T_i$（$i\in I$）的和时，6.2.1对 $\overline G$ 在 $S$ 上的\Term{section}{截面}的描述可以简化：这样的\Term{section}{截面}由一族以同一\Term{set}{集合} $I$ 编号的整数 $(n_i)_{i\in I}$，以及一族\Term{section}{截面}
\[
\overline g_i\in\Gamma(S_i',F^{(n_i)}|S_i'),\qquad i\in I,
\qquad S_i'=S-T_i',\quad T_i'=\coprod_{\substack{j\in I\\j\ne i}}T_j
\]
确定；这些 $\overline g_i$ 满足如下相容条件：对任意 $i,j\in I$，借助由这些\Term{torsor}{挠子}的\Term{section}{截面} $s^{(n_i)}$、$s^{(n_j)}$ \Term{derive}{导出}的 $F^{(n_i)}|U$、$F^{(n_j)}|U$ 的等同，$\overline g_i$ 与 $\overline g_j$ 在 $S_i'\cap S_j'=U$ 上的\Term{restriction}{限制}相同。

\textbf{6.2.3.} 最后，若假定 $T$ \Term{nonempty}{非空}且\Term{connected}{连通}，则给出 $\overline G$ 在 $S$ 上的一个\Term{section}{截面} $\overline g$，就\Term{equivalence}{等价}于给出一个整数 $n$ 以及 $F^{(n)}$ 在 $S$ 上的一个\Term{section}{截面} $\overline g$。这时不再需要相容条件，因而在描述 $\overline G$ 在 $S$ 上的\Term{section}{截面}\Term{group}{群}时不出现 $F|U$ 的\Term{section}{截面} $s$。（当然，一旦要描述 $\overline G$ 的取值于\Term{generic}{一般} $S$-\Term{scheme}{概形} $S'$ 的点，$s$ 就会再次出现。）

\textbf{6.2.4.} 由于 $\overline G$ 是6.2.1所描述的各个“\Term{open-subset}{开集}”$\alpha_n^*(\overline G)$（$n\in\mathbf Z$）的并，可知 $\overline G$ \Term{representable}{可表}当且仅当各\Term{torsor}{挠子} $F^{(n)}$ \Term{representable}{可表}。例如，由\Term{descent}{下降}理论 SGA 1，第VIII讲2.1，只要 $G$ \Term{representable}{可表}且为 $S$ 上的\Term{affine}{仿射}\Term{group-scheme}{群概形}，这个条件就成立。
\SourcePage{300}
\textbf{6.3.} 以下假定 $G=\mathbf G_{m,S}$。这样，$\mathbf Z_{S,T}$ 的每个以 $G=\mathbf G_{m,S}$ 为\Term{kernel}{核}的\Term{group-extension}{扩张} $\overline G$ 都\Term{representable}{可表}（6.2.4）；而且，作为\Term{etale}{平展}\Term{group-scheme}{群概形} $\mathbf Z_{S,T}$ 的以\Term{smooth}{光滑}\Term{group-scheme}{群概形} $\mathbf G_{m,S}$ 为\Term{kernel}{核}的\Term{group-extension}{扩张}，它当然在 $S$ 上\Term{smooth}{光滑}。此外，给出一个 $\mathbf G_m$-\Term{torsor}{挠子} $F$，\Term{equivalence}{等价}于给出 $S$ 上的一个\Term{invertible-module}{可逆模} $\mathcal L$（把后者对应到 $\mathcal L$ 的\Term{invertible}{可逆}\Term{section}{截面}所成的\Term{torsor}{挠子}）。因此，$\overline G$ 由二元组 $(\mathcal L,s)$ 描述，其中 $\mathcal L$ 是 $S$ 上的\Term{invertible-module}{可逆模}，$s$ 是 $\mathcal L$ 在 $U=S-T$ 上的一个\Term{trivialization}{平凡化}，即 $\mathcal L|U$ 的一个\Term{invertible}{可逆}\Term{section}{截面}。当 $U$ 在 $S$ 中\Term{schematically-dense}{概形稠密}，并且 $S$ \Term{locally-noetherian}{局部诺特}，或者 $S$ \Term{reduced}{约化}且其\Term{irreducible-component}{不可约分支}的\Term{set}{集合}\Term{locally-finite}{局部有限}时，给出这样的二元组，\Term{equivalence}{等价}于给出 $S$ 上一个满足
\[
\operatorname{supp}D\subset T\tag{6.3.1}
\]
的\Term{divisor}{除子} $D$；具体地，把这样的\Term{divisor}{除子}对应到二元组
\[
\bigl(\mathcal O_S(-D),s_D|U\bigr),\tag{6.3.2}
\]
其中 $\mathcal O_S(-D)$ 是 $-D$ 所关联的\Term{invertible-module}{可逆模}（EGA IV，21.2.8.1），第二项则是其\Term{canonical}{典范}\Term{rational-section}{有理截面} $s_{-D}$ 在 $U$ 上的\Term{restriction}{限制}（EGA IV，21.2.11与20.2.11（ii））。无论如何，不对 $S$ 加任何假设，$S$ 上一个满足（6.3.1）的\Term{divisor}{除子} $D$ 总能通过（6.3.2）定义二元组 $(\mathcal L,s)$，从而定义 $\mathbf Z_{S,T}$ 的以 $\mathbf G_m$ 为\Term{kernel}{核}的\Term{group-extension}{扩张} $\overline G$；我们称之为与\Term{divisor}{除子} $D$ 及满足（6.3.1）的 $S$ 的\Term{closed-subset}{闭集} $T$ 关联的 $\mathbf G_m$ 的\Term{neron-raynaud-model}{Néron--Raynaud 模型}。若 $T=\operatorname{supp}D$，则在称谓中省略 $T$。
\SourcePage{301}
显然，只要\Term{inverse-image}{逆像}\Term{divisor}{除子} $f^*(D)$ 有定义（EGA IV，21.4.2），这个\Term{group-extension}{扩张} $\overline G$ 的形成就与\Term{base-change}{基变换} $f:S'\to S$ 交换；特别地，它与任意\Term{flat}{平坦}\Term{base-change}{基变换}交换。

\textbf{6.3.3.} 由6.2.1对 $\overline G$ 的\Term{section}{截面}的\Term{generic}{一般}描述，还立即看出，这些\Term{section}{截面}可以与 $S$ 上\Term{regular}{正则}\Term{meromorphic}{亚纯}\Term{section}{截面}的\Term{sheaf}{层} $\mathcal M_S^*$（EGA IV，20.1.8）中满足如下条件的\Term{section}{截面} $f$ 等同：其\Term{divisor}{除子} $\operatorname{div}(f)$ \Term{local-scheme}{局部}地（对于 Zariski \Term{topology}{拓扑}）是 $D$ 的整数倍。若 $V$ 是 $S$ 的一个\Term{open-subset}{开集}，且 $\operatorname{div}f$ 在 $V$ 上的\Term{restriction}{限制}等于 $nD$ 在 $V$ 上的\Term{restriction}{限制}（$n\in\mathbf Z$），那么由 $f$ 定义的 $\mathbf Z_{S,T}$ 的\Term{section}{截面}，正是（6.1.2）中 $\mathbf Z_S$ 的值为 $n$ 的\Term{constant}{常值}\Term{section}{截面}的像。上述对 $\overline G$ 在 $S$ 上的\Term{section}{截面}\Term{group}{群}的描述，在 $S$ 的任意\Term{open-subset}{开集}上仍成立；更\Term{generic}{一般}地，对 $\overline G$ 的取值于任意\Term{flat}{平坦} $S$-\Term{scheme}{概形} $S'$ 的点所成的\Term{group}{群}也成立。由于 $\overline G$ 本身在 $S$ 上\Term{flat}{平坦}（甚至\Term{smooth}{光滑}），这个描述便再次完整刻画了由 $D$ 及 $T\supset\operatorname{supp}D$ 定义的 $\mathbf G_m$ 的\Term{neron-raynaud-model}{Néron--Raynaud 模型}。

\textbf{6.4.} 现在假定 $S$ \Term{normal}{正规}，且其\Term{irreducible-component}{不可约分支}的\Term{set}{集合}\Term{locally-finite}{局部有限}（于是 $S$ 是\Term{normal}{正规}\Term{integral-scheme}{整概形}的和），并假定\Term{divisor}{除子} $D$ 为\Term{effective-divisor}{正除子}，各\Term{multiplicity}{重数}均为 $1$（EGA IV，21.6），$T=\operatorname{supp}D$，且 $T$ \Term{locally-irreducible}{局部不可约}，即其\Term{irreducible-component}{不可约分支}两两不交。于是，对 $S$ 上每个在 $U=S-T$ 上\Term{regular}{正则}的\Term{rational-function}{有理函数} $f$，$\operatorname{div}f$ \Term{local-scheme}{局部}地是 $D$ 的倍数；因此，这时有\Term{canonical}{典范}的\Term{group}{群}\Term{isomorphism}{同构}
\SourcePage{302}
\[
\Gamma(S,\overline G)\xrightarrow{\sim}\Gamma(U,\mathbf G_{m,U}).\tag{6.4.1}
\]
进一步假定 $T$ \Term{geometrically-unibranch}{几何单分支}且\Term{locally-integral}{局部整}；这个假设在\Term{smooth}{光滑}\Term{base-change}{基变换}下保持\Term{invariant}{不变}（EGA IV，17.5.7）。于是，由于所有假设都在\Term{base-change}{基变换}下保持\Term{invariant}{不变}，由（6.4.1）可知，对 $S$ 上\Term{smooth}{光滑}的 $S'$，有关于 $S'$ \Term{functoriality}{函子性}的\Term{canonical}{典范}\Term{group}{群}\Term{isomorphism}{同构}
\[
\operatorname{Hom}_S(S',\overline G)\xrightarrow{\sim}
\operatorname{Hom}_U(S_U',\mathbf G_{m,U})\simeq\Gamma(S_U',\mathcal O_{S_U'}^*).\tag{6.4.2}
\]
由于 $\overline G$ 本身在 $S$ 上\Term{smooth}{光滑}，这个\Term{isomorphism}{同构}在唯一\Term{isomorphism}{同构}意义下刻画了 $S$ 上的\Term{group-scheme}{群概形} $\overline G$。另一方面，它表明 $\overline G$ 对 $S$ 的\Term{open-subset}{开集} $U$ 上的\Term{group-scheme}{群概形} $\mathbf G_{m,U}$ 起到“\Term{neron-model}{Néron 模型}”的\Term{action}{作用}；再加上它最早由 M. Raynaud 引入并使用，这就说明了6.3中用以称呼 $\overline G$ 的术语的来由。
\begin{SGAresult}{\Term{theorem}{定理}}{6.5}
设 $S$ 为\Term{normal}{正规}且\Term{locally-integral}{局部整}的\Term{scheme}{概形}，$D$ 为 $S$ 上没有多重\Term{branch}{分支}的\Term{effective-divisor}{正除子}，其\Term{support}{支集}\Term{geometrically-unibranch}{几何单分支}且\Term{locally-integral}{局部整}。设 $\overline G$ 为 $S$ 上关于 $D$ 的 $\mathbf G_m$ 的\Term{neron-raynaud-model}{Néron--Raynaud 模型}（6.3），考虑 $S$ 上以\Term{group-scheme}{群概形} $\overline G$ 为结构\Term{group}{群}的\Term{torsor}{挠子} $E$ 的\Term{category}{范畴}。\Term{restriction}{限制}\Term{functor}{函子}
\[
E\longmapsto E_U=E|U
\]
从这个\Term{category}{范畴}到 $U$ 上以\Term{group-scheme}{群概形} $\mathbf G_{m,U}$ 为结构\Term{group}{群}的\Term{torsor}{挠子}的\Term{category}{范畴}，是\Term{fully-faithful}{充分忠实}的。若 $S$ 是\Term{regular}{正则}\Term{scheme}{概形}（更\Term{generic}{一般}地，若其\Term{local-ring}{局部环}都\Term{factorial}{析因}），则这个\Term{functor}{函子}甚至是\Term{equivalence-of-categories}{范畴等价}。
\end{SGAresult}
\SourcePage{303}
\begin{proof}
先证明所考虑的\Term{functor}{函子}\Term{fully-faithful}{充分忠实}，即对两个 $\overline G$-\Term{torsor}{挠子} $E,F$，\Term{restriction-map}{限制映射}
\[
\operatorname{Hom}(E,F)\longrightarrow\operatorname{Hom}(E_U,F_U)\tag{*}
\]
是\Term{bijection}{双射}，这里的 $\operatorname{Hom}$ 表示\Term{torsor}{挠子}\Term{homomorphism}{同态}的\Term{set}{集合}。首先注意，这个问题显然对于 $S$ 上的 Zariski \Term{topology}{拓扑}是\Term{local-scheme}{局部}的：若 $S$ 是一些\Term{open-subset}{开集} $S_i$ 的并，且对每个包含于某个 $S_i$ 的 $S$ 的\Term{open-subset}{开集} $S'$（特别是 $S_i$ 或 $S_i\cap S_j$ 这样的 $S'$），映射
\[
\operatorname{Hom}(E_{S'},F_{S'})\longrightarrow
\operatorname{Hom}(E_{S'\cap U},F_{S'\cap U})
\]
都是\Term{bijection}{双射}，那么（*）也是\Term{bijection}{双射}。另一方面，把5.8应用于 $S,T,\overline G$，把5.1应用于 $T,\mathbf Z$，便立即可知每个 $\overline G$-\Term{torsor}{挠子}对于 Zariski \Term{topology}{拓扑}都\Term{locally-trivial}{局部平凡}。这两点把问题归结为在 $E=F=\overline G$ 时证明（*）为\Term{bijection}{双射}。但此时这个映射就是（6.4.1），它的\Term{bijection}{双射}性已经证明。

现在假定 $X$ 的\Term{local-ring}{局部环}都\Term{factorial}{析因}，证明所考虑的\Term{functor}{函子} $E\mapsto E_U$ \Term{essentially-surjective}{本质满射}。由于 $\mathbf G_{m,S}$ 等同于 $\overline G$ 的一个\Term{open-subgroup-scheme}{开子群概形}，只须证明每个 $\mathbf G_{m,U}$-\Term{torsor}{挠子}都\Term{isomorphism}{同构}于某个 $\mathbf G_{m,S}$-\Term{torsor}{挠子}在 $U$ 上的\Term{restriction}{限制}。这\Term{equivalence}{等价}于说，$U$ 上的每个\Term{invertible-module}{可逆模}都\Term{isomorphism}{同构}于 $S$ 上某个\Term{invertible-module}{可逆模}的\Term{restriction}{限制}，而这确实由对 $S$ 所作的假设得到（EGA IV，21.6.11；该处限于\Term{locally-noetherian}{局部诺特}的情形，这也是我们以后唯一会用到的情形）。这就证明了6.5。
\end{proof}
\SourcePage{304}
\textbf{6.5.1.} 现在进一步假定 $S$ \Term{locally-noetherian}{局部诺特}且\Term{regular}{正则}。由于 $S$ 此时\Term{locally-factorial}{局部析因}，给出 $D$，\Term{equivalence}{等价}于给出一个\Term{pure-codimension-one}{纯余维为 $1$}、\Term{reduced}{约化}且\Term{geometrically-unibranch}{几何单分支}的\Term{closed-subscheme}{闭子概形}（仍记为 $D$）。还应注意，对 $S,D$ 的假设在任意\Term{smooth}{光滑}\Term{base-change}{基变换} $S'\to S$ 下保持\Term{invariant}{不变}，所以6.5的结论适用于配备\Term{closed-subscheme}{闭子概形} $D'=D\times_S S'$ 的 $S'$。于是，第VII讲1.1.6对\Term{group-extension}{群扩张}的描述立即给出下述结果。
\begin{SGAresult}{推论}{6.6}
设 $S$ 为\Term{locally-noetherian}{局部诺特}且\Term{regular}{正则}的\Term{scheme}{概形}，$D$ 为\Term{reduced}{约化}、\Term{geometrically-unibranch}{几何单分支}且\Term{pure-codimension-one}{纯余维为 $1$}的\Term{closed-subscheme}{闭子概形}（因而定义 $S$ 上的一个\Term{divisor}{除子}），$\overline G$ 为关于 $D$ 的 $\mathbf G_m$ 的\Term{neron-raynaud-model}{Néron--Raynaud 模型}，$P$ 为 $S$ 上\Term{smooth}{光滑}的\Term{commutative-group-scheme}{交换群概形}。考虑\Term{group-sheaf}{群层} $P$ 的以 $\overline G$ 为\Term{kernel}{核}的\Term{commutative-group-extension}{交换扩张}的\Term{category}{范畴}。则\Term{restriction}{限制}到 $U$ 的\Term{functor}{函子} $E\mapsto E|U$ 在这个\Term{category}{范畴}与 $P_U$ 的以 $\mathbf G_{m,U}$ 为\Term{kernel}{核}的\Term{group-extension}{扩张}的\Term{category}{范畴}之间给出\Term{equivalence-of-categories}{范畴等价}。
\end{SGAresult}
\textbf{6.6.1.} 同一方法给出\Term{group-extension}{扩张}不必交换时的类似命题（此时也不再假定 $P$ 交换）。证明的细节留给读者；由于这里要考虑\Term{bitorsor}{双挠子}而不是\Term{torsor}{挠子}，证明稍复杂一些。

用与6.6本质相同的证明，还得到：
\SourcePage{305}
\begin{SGAresult}{推论}{6.7}
设 $S,D,\overline G$ 如6.6，$P,Q$ 为 $S$ 上的两个\Term{smooth}{光滑}\Term{commutative-group-scheme}{交换群概形}。考虑 $(P,Q)$ 的以 $\overline G$ 为\Term{kernel}{核}的\Term{biextension}{双扩张}的\Term{category}{范畴}。则\Term{restriction}{限制}到 $U$ 的\Term{functor}{函子} $E\mapsto E|U$，从这个\Term{category}{范畴}到 $(P_U,Q_U)$ 的以 $\mathbf G_{m,U}$ 为\Term{kernel}{核}的\Term{biextension}{双扩张}的\Term{category}{范畴}，是\Term{equivalence-of-categories}{范畴等价}。
\end{SGAresult}
\textbf{注记6.8.} a）我们将在 $S$ 为\Term{trait}{离散赋值环谱}、$D$ 为由 $S$ 的\Term{closed-point}{闭点} $s$ 到 $S$ 的包含定义的\Term{divisor}{除子}时应用6.6与6.7。这时 $U$ 是 $S$ 的\Term{fraction-field}{分式域} $K$ 的\Term{spectrum}{谱}，$\overline G$ 是 Raynaud 引入的 $K$ 上\Term{multiplicative-group}{乘法群}的通常\Term{neron-model}{Néron 模型}。应当注意，即使仅限于这个情形，6.6与6.7的证明仍使用了比\Term{trait}{离散赋值环谱}更\Term{generic}{一般}的底（如 $P$ 及其各\Term{cartesian-power}{笛卡儿幂}等）上的6.5。然而，在这种情形，只须知道所考虑的特殊情形下\Term{neron-model}{Néron 模型} $\overline G$ 的存在性，就能\Term{form}{形式}地\Term{pushout}{推出}在 $S$ 上\Term{smooth}{光滑}的底\Term{scheme}{概形}上的6.5，进而得到6.6与6.7。对我们而言，引入 $\overline G$ 主要是为了给下一节研究以\Term{multiplicative-group}{乘法群}为\Term{kernel}{核}的\Term{group-extension}{扩张}与\Term{biextension}{双扩张}的延拓问题提供一个方便的中介。

 b）在远比以上考虑的条件更\Term{generic}{一般}的条件下，也可以定义 $\mathbf G_{m,U}$ 的\Term{neron-raynaud-model}{Néron--Raynaud 模型} $\overline G$，使其对 $S$ 上每个\Term{smooth}{光滑}的 $S'$ 都满足基本条件（6.4.2）：这里只要求 $U$ 是\Term{locally-noetherian}{局部诺特}\Term{normal}{正规}\Term{scheme}{概形} $S$ 的一个\Term{schematically-dense}{概形稠密}\Term{open-subset}{开集}。此时 $\overline G$ 是一个 fppf \Term{sheaf}{层}，不一定\Term{representable}{可表}。用 $\mathcal D$ 表示 $S$ 的小\Term{etale}{平展}\Term{site}{景}上，$S$ 上\Term{support}{支集}包含于 $T=S-U$ 的\Term{divisor}{除子}所成的\Term{sheaf}{层}，并略用同一个字母表示它在 $(\mathrm{Sch})/S$ 上的\Term{canonical}{典范}延伸\Term{sheaf}{层}（其在相对\Term{scheme}{概形} $f:S'\to S$ 上的\Term{section}{截面}是 $f_{\text{ét}}^*(\mathcal D)$ 在 $S'$ 上的\Term{section}{截面}），则所求\Term{model}{模型}是一个\Term{group-extension}{群扩张}
\SourcePage{306}
\[
0\longrightarrow\mathbf G_{m,S}\longrightarrow\overline G\longrightarrow\mathcal D\longrightarrow0.\tag{6.8.1}
\]
利用第VII讲1.1.6的\Term{generic}{一般}描述，构造一个\Term{canonical}{典范}\Term{group-extension}{扩张}（6.8.1）并无困难（无需对 $S$ 作\Term{normal}{正规}性假设），同时还可得到 $S'$ 在 $S$ 上\Term{etale}{平展}时的（6.4.2）。要使同一关系对 $S$ 上任意\Term{smooth}{光滑}的 $S'$ 成立，\Term{equivalence}{等价}于验证\Term{sheaf}{层} $\mathcal D$ 的形成与任意\Term{smooth}{光滑}\Term{base-change}{基变换} $f:S'\to S$ 交换；当 $S$ \Term{normal}{正规}时，这由 EGA $\mathrm{Err}_{\mathrm{IV}}$ 53，21.4.9得到。

\textbf{6.9.“自我批评”（P. Deligne 的意见）}\footnote{本讲的执笔者承认无能为力；若有需要，请读者依照附在这里的指示重写第5、6节。}

第5、6节实在不大能启发人。读着读着，总觉得是在“手工”重做 $i_*,i^*$ 的\Term{form}{形式}体系，却看不出哪些内容是几何的，哪些只是“\Term{generic}{一般}废话”（general non-sense）。

我建议提炼出下面这些“几何”命题 a）、b）、c）：

 a）（参见6.8 b）设 $j:U\hookrightarrow S$ 是\Term{normal}{正规}\Term{noetherian}{诺特}\Term{scheme}{概形} $S$ 的一个\Term{schematically-dense}{概形稠密}\Term{open-subset}{开集}。于是\Term{etale}{平展}\Term{sheaf}{层} $\mathbf G_{m,S}$ 等同于\Term{etale}{平展}\Term{sheaf}{层} $j_*\mathbf G_{m,U}$ 的一个\Term{subsheaf}{子层}。设 $\mathcal D$ 为其\Term{quotient-sheaf}{商层}：
\[
0\longrightarrow\mathbf G_{m,S}\longrightarrow j_*\mathbf G_{m,U}\longrightarrow\mathcal D\longrightarrow0.
\]
\SourcePage{307}
则 $\mathcal D$ 的形成与任意\Term{smooth}{光滑}\Term{base-change}{基变换}交换（EGA $\mathrm{Err}_{\mathrm{IV}}$ 53，21.4.9）；此外，若 $S$ 的\Term{strictly-local-ring}{严格局部环}都\Term{factorial}{析因}，则 $R^1j_*\mathbf G_m=0$（\Term{theorem}{定理}90）。

 b）若 $U=S-Y$，其中 $Y$ 为一个\Term{geometrically-unibranch}{几何单分支}\Term{divisor}{除子}，且 $i$ 为 $Y$ 到 $S$ 的包含，则
\[
\mathcal D=i_*\mathbf Z_Y.
\]
 c）若 $Y$ 是\Term{scheme}{概形} $X$ 的一个\Term{closed-subscheme}{闭子概形}，$E$ 是\Term{set}{集合}，则\Term{etale}{平展}\Term{sheaf}{层} $i_*E_Y$ 由 $X$ 上的一个\Term{etale}{平展}\Term{scheme}{概形}表示，其形成与任意\Term{base-change}{基变换}相容。
\begin{SGAresult}{推论}{}
设 $S$ 为\Term{noetherian}{诺特}\Term{normal}{正规}\Term{scheme}{概形}，$Y$ 为一个\Term{geometrically-unibranch}{几何单分支}\Term{divisor}{除子}，$U=S-Y$：
\[
\begin{tikzcd}
U\arrow[r,hook,"j"]&S&Y\arrow[l,hook',"i"']
\end{tikzcd}
\]
存在唯一的 $S$ 上\Term{smooth}{光滑}\Term{group-scheme}{群概形} $\mathscr G_m$，使得在任意\Term{smooth}{光滑}\Term{base-change}{基变换}后，都有
\[
\mathscr G_m|S_{\text{ét}}=j_*(\mathbf G_m)_U.
\]
\end{SGAresult}
\begin{proof}
上述要求确实确定了 $\mathscr G_m$ 作为 $S$ 上\Term{smooth}{光滑}\Term{scheme}{概形}所成\Term{category}{范畴}上的\Term{sheaf}{层}（对于\Term{etale}{平展}\Term{topology}{拓扑}）的结构。由 a）、b），序列
\[
0\longrightarrow\mathbf G_m\longrightarrow\mathscr G_m\xrightarrow{q}i_*\mathbf Z_Y\longrightarrow0
\]
是\Term{exact-sequence}{正合列}。由 c），$i_*\mathbf Z_Y$ \Term{representable}{可表}。最后，\Term{morphism}{态射} $q$ 也\Term{representable}{可表}。
\end{proof}
使用上述“小\Term{smooth}{光滑}\Term{site}{景}”，应当能够从 a）\Term{form}{形式}地\Term{pushout}{推出}第6节的各个命题。
\SourcePage{308}
\SGAsection{7}{以乘法群为核的扩张与双扩张的典范延拓}
\textbf{7.0.} 现在可以收获第5、6节准备工作的成果了。先设底\Term{scheme}{概形} $S$ 是\Term{trait}{离散赋值环谱}，即一个\Term{discrete-valuation-ring}{离散赋值环}的\Term{spectrum}{谱}，其\Term{generic-point}{泛点}为 $\eta$，\Term{closed-point}{闭点}为 $s$，$k=k(s)$。设 $P$ 为 $S$ 上\Term{smooth}{光滑}的\Term{commutative-group-scheme}{交换群概形}，记 $P_0=P_s$，即 $P$ 的\Term{special-fiber}{特殊纤维}。本节只考虑\Term{commutative-group}{交换群}，特别地，所考虑的\Term{group-extension}{群扩张}都是交换的。
\begin{SGAresult}{\Term{theorem}{定理}}{7.1}
记号如7.0。

 a）假定 $P$ 的\Term{special-fiber}{特殊纤维}的\Term{connected-component-group}{连通分支群}
\[
\Phi=P_0/P_0^0
\]
是\Term{torsion-group}{挠群}；特别地，当 $P$ 在 $S$ 上\Term{finite-type}{有限型}时就是如此。则\Term{functor}{函子} $E\mapsto E_\eta$，从 $P$ 的以 $\mathbf G_{m,S}$ 为\Term{kernel}{核}的\Term{group-extension}{扩张}的\Term{category}{范畴}到 $P_\eta$ 的以 $\mathbf G_{m,\eta}$ 为\Term{kernel}{核}的\Term{group-extension}{扩张}的\Term{category}{范畴}，是\Term{fully-faithful}{充分忠实}的。假定存在整数 $n>0$ 使 $n\Phi=0$，则对 $P_\eta$ 的以 $\mathbf G_{m,\eta}$ 为\Term{kernel}{核}的\Term{group-extension}{扩张} $E_\eta$，有一个与之关联的\Term{canonical}{典范}\Term{homomorphism}{同态}
\[
d(E_\eta):\Phi\longrightarrow(\mathbf Q/\mathbf Z)_k,\tag{7.1.1}
\]
它为零，是\Term{group-extension}{扩张} $E_\eta$ \Term{isomorphism}{同构}于 $P$ 的某个以 $\mathbf G_{m,S}$ 为\Term{kernel}{核}的\Term{group-extension}{扩张}的\Term{generic-fiber}{泛纤维}的充分必要条件。特别地，若 $P_0$ \Term{connected}{连通}，则上述\Term{functor}{函子} $E\mapsto E_\eta$ 是\Term{equivalence-of-categories}{范畴等价}。
\SourcePage{309}
 b）设 $P,Q$ 为 $S$ 上两个\Term{smooth}{光滑}\Term{group-scheme}{群概形}，置
\[
\Phi=P_0/P_0^0,\qquad\Psi=Q_0/Q_0^0,
\]
并假定 $\Phi$ 或 $\Psi$ 是\Term{torsion-group}{挠群}（例如 $P$ 或 $Q$ 在 $S$ 上\Term{finite-type}{有限型}）；则\Term{functor}{函子} $E\mapsto E_\eta$，从 $(P,Q)$ 的以 $\mathbf G_{m,S}$ 为\Term{kernel}{核}的\Term{biextension}{双扩张}的\Term{category}{范畴}到 $(P_\eta,Q_\eta)$ 的以 $\mathbf G_{m,\eta}$ 为\Term{kernel}{核}的\Term{biextension}{双扩张}的\Term{category}{范畴}，是\Term{fully-faithful}{充分忠实}的。假定存在整数 $n>0$ 使 $n\Phi=0$ 或 $n\Psi=0$，则对 $(P_\eta,Q_\eta)$ 的以 $\mathbf G_{m,\eta}$ 为\Term{kernel}{核}的\Term{biextension}{双扩张} $E_\eta$，有一个与之关联的\Term{canonical}{典范}\Term{pairing}{配对}
\[
d(E_\eta):\Phi\otimes_{\mathbf Z}\Psi\longrightarrow(\mathbf Q/\mathbf Z)_k,\tag{7.1.2}
\]
它为零，是 $E_\eta$ \Term{isomorphism}{同构}于 $(P,Q)$ 的某个以 $\mathbf G_{m,S}$ 为\Term{kernel}{核}的\Term{biextension}{双扩张} $E$ 的\Term{generic-fiber}{泛纤维}的充分必要条件。特别地，若 $P_0$ 或 $Q_0$ \Term{connected}{连通}，则上述\Term{functor}{函子} $E\mapsto E_\eta$ 是\Term{equivalence-of-categories}{范畴等价}。
\end{SGAresult}
\begin{proof}
 a）设 $\overline G$ 为 $\mathbf G_m$ 的\Term{neron-raynaud-model}{Néron--Raynaud 模型}。所考虑的\Term{functor}{函子}是“\Term{kernel}{核}的扩充 $\mathbf G_{m,S}\to\overline G$”\Term{functor}{函子}与定义在以 $\overline G$ 为\Term{kernel}{核}的\Term{biextension}{双扩张}的\Term{category}{范畴}上的“\Term{restriction}{限制}到 $\eta$”\Term{functor}{函子}的复合。由6.6，后一个\Term{functor}{函子}是\Term{equivalence-of-categories}{范畴等价}，故只须研究第一个\Term{functor}{函子}。为此，使用\Term{group-scheme}{群概形}的\Term{exact-sequence}{正合列}
\[
0\longrightarrow\mathbf G_{m,S}\longrightarrow\overline G\longrightarrow\mathbf Z_{S,s}\longrightarrow0,\tag{7.1.3}
\]
其中 $\mathbf Z_{S,s}$ 的记号如5.7。与这个\Term{exact-sequence}{正合列}关联的 $\operatorname{Ext}^i$（$i=0,1$）\Term{exact-sequence}{正合列}表明，所考虑的\Term{functor}{函子}\Term{fully-faithful}{充分忠实}的充分必要条件是
\SourcePage{310}
\[
\operatorname{Hom}(P,\mathbf Z_{S,s})=0,\tag{7.1.4}
\]
而 $P$ 的一个以 $\overline G$ 为\Term{kernel}{核}的\Term{group-extension}{扩张} $\overline E$ 来自以 $\mathbf G_{m,S}$ 为\Term{kernel}{核}的\Term{group-extension}{扩张} $E$ 的\Term{obstruction}{障碍}位于
\[
\operatorname{Ext}^1(P,\mathbf Z_{S,s}).\tag{7.1.5}
\]
由（5.7.2），（7.1.4）中的第一个\Term{group}{群}\Term{isomorphism}{同构}于 $\operatorname{Hom}(P_0,\mathbf Z_k)$，后者又\Term{isomorphism}{同构}于 $\operatorname{Hom}(\Phi,\mathbf Z_k)$；所以（7.1.4）\Term{equivalence}{等价}于关系
\[
\operatorname{Hom}(\Phi,\mathbf Z_k)=0.\tag{7.1.4 bis}
\]
当 $\Phi$ 为\Term{torsion-group}{挠群}时，这个关系当然成立。另一方面，若更有 $n\Phi=0$，则由5.9与5.5，有\Term{canonical-isomorphism}{典范同构}
\[
\operatorname{Ext}^1(P,\mathbf Z_{S,s})\simeq
\operatorname{Ext}^1(P_0,\mathbf Z_k)\simeq
\operatorname{Hom}\bigl(\Phi,(\mathbf Q/\mathbf Z)_k\bigr),\tag{7.1.6}
\]
这就完成了7.1 a）的证明。

 b）证明本质相同。由6.7，只须研究“\Term{kernel}{核}的扩充 $\mathbf G_{m,S}\to\overline G$”这个\Term{functor}{函子}：它从 $(P,Q)$ 的以 $\mathbf G_{m,S}$ 为\Term{kernel}{核}的\Term{biextension}{双扩张}的\Term{category}{范畴}到 $(P,Q)$ 的以 $\overline G$ 为\Term{kernel}{核}的\Term{biextension}{双扩张}的\Term{category}{范畴}。使用\Term{exact-sequence}{正合列}（7.1.3）及其关联的 $\operatorname{Biext}^i$（$i=0,1$）\Term{exact-sequence}{正合列}（第VII讲3.7.5），即可进行这一研究。所考虑的\Term{functor}{函子}\Term{fully-faithful}{充分忠实}\Term{equivalence}{等价}于关系
\[
\operatorname{Hom}(P\otimes Q,\mathbf Z_{S,s})=0.\tag{7.1.7}
\]
由（5.7.2），左端\Term{isomorphism}{同构}于 $\operatorname{Hom}(P_0\otimes Q_0,\mathbf Z_k)$，因而\Term{isomorphism}{同构}于 $\operatorname{Hom}(\Phi\otimes\Psi,\mathbf Z_k)$；若 $\Phi$ 或 $\Psi$ 是\Term{torsion-group}{挠群}，这个\Term{group}{群}确实为零。
\SourcePage{311}
另一方面，对 $(P,Q)$ 的给定的以 $\overline G$ 为\Term{kernel}{核}的\Term{biextension}{双扩张} $\overline E$，它来自 $(P,Q)$ 的以 $\mathbf G_{m,S}$ 为\Term{kernel}{核}的\Term{biextension}{双扩张} $E$ 的\Term{obstruction}{障碍}位于 $\operatorname{Biext}^1(P,Q;\mathbf Z_{S,s})$；由5.10，这\Term{isomorphism}{同构}于 $\operatorname{Biext}^1(P_0,Q_0;\mathbf Z_k)$；当 $n\Phi=0$ 或 $n\Psi=0$ 时，后者又由5.6\Term{isomorphism}{同构}于 $\operatorname{Hom}(\Phi\otimes\Psi,\mathbf Q/\mathbf Z)$。这就完成了7.1的证明。
\end{proof}
\textbf{注记7.2.} 给定\Term{generic-point}{泛点}为 $\eta$ 的\Term{locally-noetherian}{局部诺特}\Term{regular}{正则}\Term{connected}{连通}\Term{scheme}{概形} $S$，以及 $S$ 上的\Term{smooth}{光滑}\Term{group-scheme}{群概形} $P$，可以研究从 $P$ 的以 $\mathbf G_{m,S}$ 为\Term{kernel}{核}的\Term{group-extension}{扩张}的\Term{category}{范畴}到 $P_\eta$ 的以 $\mathbf G_{m,\eta}$ 为\Term{kernel}{核}的\Term{group-extension}{扩张}的\Term{category}{范畴}的\Term{functor}{函子} $E\mapsto E_\eta$。事实上，这个问题容易归结为7.1 a）处理的问题。为简便起见，假定 $P$ 在 $S$ 上\Term{finite-type}{有限型}。由7.1容易得到，所考虑的\Term{functor}{函子}\Term{fully-faithful}{充分忠实}\footnote{为此，只须 $S$ \Term{normal}{正规}，不必\Term{regular}{正则}。}。由此容易\Term{pushout}{推出}，若 $E_\eta$ 是 $P_\eta$ 的以 $\mathbf G_{m,\eta}$ 为\Term{kernel}{核}的\Term{group-extension}{扩张}，则它\Term{isomorphism}{同构}于 $P$ 的某个以 $\mathbf G_{m,S}$ 为\Term{kernel}{核}的\Term{group-extension}{扩张} $E$ 的\Term{generic-fiber}{泛纤维}，当且仅当对每个满足 $\dim\operatorname{Spec}(\mathcal O_{S,s})=1$ 的 $s\in S$，由\Term{base-change}{基变换} $\operatorname{Spec}(\mathcal O_{S,s})\to S$ 得到的类似延拓问题有解；由7.1 a），这可表述为某个\Term{homomorphism}{同态}
\[
d(E_\eta,s):P_s/P_s^0\longrightarrow(\mathbf Q/\mathbf Z)_{k(s)}\tag{7.2.1}
\]
为零。上述断言显然归结为下述命题：若 $U$ 是 $S$ 的一个\Term{open-subset}{开集}，其补集 $T$ 满足
\[
\operatorname{codim}(T,S)\geq2,
\]
\SourcePage{312}
则\Term{functor}{函子} $E\mapsto E|U$，从 $P$ 的以 $\mathbf G_{m,S}$ 为\Term{kernel}{核}的\Term{group-extension}{扩张}的\Term{category}{范畴}到 $P_U$ 的以 $\mathbf G_{m,U}$ 为\Term{kernel}{核}的\Term{group-extension}{扩张}的\Term{category}{范畴}，是\Term{equivalence-of-categories}{范畴等价}。而从第VII讲1.1.6的描述来看，这个断言本质上是显然的；这里使用如下事实：对任意\Term{smooth}{光滑} $S$-\Term{scheme}{概形} $X$（如 $P,P\times P,P\times P\times P,\ldots$），$X$ \Term{locally-factorial}{局部析因}，且 $X_T=X-X_U$ 在 $X$ 中\Term{codimension}{余维}至少为 $2$，从而\Term{functor}{函子} $F\mapsto F|X_U$，从 $\mathbf G_{m,X}$-\Term{torsor}{挠子}的\Term{category}{范畴}到 $\mathbf G_{m,X_U}$-\Term{torsor}{挠子}的\Term{category}{范畴}，是\Term{equivalence-of-categories}{范畴等价}。

关于\Term{biextension}{双扩张}的延拓问题，有完全类似的结果，具体表述留给读者。

\textbf{7.3.} 我们将具体说明（7.1.1）和（7.1.2）定义的\Term{obstruction}{障碍}在 $P,Q,S$ 变化时的若干\Term{functoriality}{函子性质}，以及这些\Term{obstruction}{障碍}彼此的关系。以下总假定定义这些\Term{obstruction}{障碍}所需的条件成立。由于验证没有困难，留给读者。

\textbf{7.3.1.} 设
\[
u:P'\longrightarrow P
\]
是 $S$ 上\Term{smooth}{光滑}\Term{group-scheme}{群概形}的一个\Term{morphism}{态射}，$E_\eta$ 为 $P_\eta$ 的以 $\mathbf G_{m,\eta}$ 为\Term{kernel}{核}的\Term{group-extension}{扩张}。沿 $u_\eta:P_\eta'\to P_\eta$ 取\Term{inverse-image}{逆像}，得到 $P_\eta'$ 的以 $\mathbf G_{m,\eta}$ 为\Term{kernel}{核}的\Term{group-extension}{扩张} $E_\eta'$。于是下图交换：
\SourcePage{313}
\[
\begin{tikzcd}[column sep=large]
P_0'/P_0'^0\arrow[d,"u_0"']\arrow[dr,"d(E_\eta')"]&\\
P_0/P_0^0\arrow[r,"d(E_\eta)"']&(\mathbf Q/\mathbf Z)_k
\end{tikzcd}\tag{7.3.1.1}
\]
同样，若有 $S$ 上\Term{smooth}{光滑}\Term{group-scheme}{群概形}的\Term{morphism}{态射}
\[
u:P'\longrightarrow P,\qquad v:Q'\longrightarrow Q,
\]
以及 $(P_\eta,Q_\eta)$ 的一个以 $\mathbf G_{m,\eta}$ 为\Term{kernel}{核}的\Term{biextension}{双扩张} $E_\eta$，沿 $(u_\eta,v_\eta)$ 取\Term{inverse-image}{逆像}得到 $(P_\eta',Q_\eta')$ 的以 $\mathbf G_{m,\eta}$ 为\Term{kernel}{核}的\Term{biextension}{双扩张} $E_\eta'$，则下图交换：
\[
\begin{tikzcd}[column sep=large]
P_0'/P_0'^0\otimes Q_0'/Q_0'^0
 \arrow[d,"u_0\otimes v_0"']\arrow[dr,"d(E_\eta')"]&\\
P_0/P_0^0\otimes Q_0/Q_0^0\arrow[r,"d(E_\eta)"']&(\mathbf Q/\mathbf Z)_k
\end{tikzcd}\tag{7.3.1.2}
\]
这些相容性的验证本质上只是机械计算，留给读者。

\textbf{7.3.2.} 由相容性（7.3.1.1）可知，若 $E_\eta$ 是 $P_\eta$ 的以 $\mathbf G_{m,\eta}$ 为\Term{kernel}{核}的\Term{group-extension}{扩张}，则存在 $P$ 的一个最大的\Term{subgroup-scheme}{子群概形} $V$，使得 $E_\eta$ 在 $V_\eta$ 上的\Term{restriction}{限制}可以延拓为 $V$ 的以 $\mathbf G_{m,S}$ 为\Term{kernel}{核}的\Term{group-extension}{扩张}：这就是 $P$ 中包含 $P_\eta$、由条件
\SourcePage{314}
\[
V_0/V_0^0=\operatorname{Ker}d(E_\eta)\tag{7.3.2.1}
\]
定义的\Term{open-subgroup}{开子群} $V$。于是，若 $u:P'\to P$ 是 $S$ 上\Term{smooth}{光滑}\Term{group-scheme}{群概形}的\Term{morphism}{态射}，则 $E_\eta$ 沿 $u_\eta$ 的\Term{inverse-image}{逆像}\Term{group-extension}{扩张} $E_\eta'$ 可延拓为 $P'$ 的以 $\mathbf G_{m,S}$ 为\Term{kernel}{核}的\Term{group-extension}{扩张}，当且仅当
\[
u(P')\subset V.\tag{7.3.2.2}
\]
对 $(P_\eta,Q_\eta)$ 的一个以 $\mathbf G_{m,\eta}$ 为\Term{kernel}{核}的\Term{biextension}{双扩张} $E_\eta$，同样存在 $P$ 的一个最大的\Term{subgroup-scheme}{子群概形} $V$，使得 $E_\eta$ 在 $(V_\eta,Q_\eta)$ 上的\Term{restriction}{限制}可以延拓为 $(V,Q)$ 的以 $\mathbf G_{m,S}$ 为\Term{kernel}{核}的\Term{biextension}{双扩张}：这就是 $P$ 中包含\Term{generic-fiber}{泛纤维} $P_\eta$、满足如下条件的\Term{open-subgroup}{开子群}：$V_0/V_0^0$ 为 $\Phi$ 中使\Term{pairing}{配对}（7.1.2）在左侧为零的\Term{subgroup}{子群}，即相应的从 $\Phi$ 到 $\underline{\operatorname{Hom}}(\Psi,(\mathbf Q/\mathbf Z)_k)$ 的\Term{homomorphism}{同态}的\Term{kernel}{核}：
\[
V_0/V_0^0=\operatorname{Ker}\left(
\Phi\xrightarrow{d(E_\eta)}\underline{\operatorname{Hom}}\bigl(\Psi,(\mathbf Q/\mathbf Z)_k\bigr)\right).\tag{7.3.2.3}
\]
于是，给定 $S$ 上\Term{smooth}{光滑}\Term{group-scheme}{群概形}的 $u:P'\to P$，由 $E_\eta$ 沿 $u_\eta$ 取\Term{inverse-image}{逆像}所得的 $(P_\eta',Q_\eta)$ 的以 $\mathbf G_{m,\eta}$ 为\Term{kernel}{核}的\Term{biextension}{双扩张}，可延拓为 $(P',Q)$ 的以 $\mathbf G_{m,S}$ 为\Term{kernel}{核}的\Term{biextension}{双扩张}，当且仅当关系（7.3.2.2）成立。交换 $P,Q$ 的地位，有\Term{symmetric}{对称}的结果。但当然，\Term{generic}{一般}找不到最大的一对\Term{subgroup}{子群} $P'\to P$、$Q'\to Q$，使 $E_\eta$ 在 $(P_\eta',Q_\eta')$ 上的\Term{restriction}{限制}\Term{biextension}{双扩张} $E_\eta'$ 可延拓为 $(P',Q')$ 的以 $\mathbf G_{m,S}$ 为\Term{kernel}{核}的\Term{biextension}{双扩张}；只能说，对给定的一对 $(P',Q')$，更\Term{generic}{一般}地，对给定的一对\Term{smooth}{光滑}\Term{group-scheme}{群概形}的\Term{morphism}{态射} $P'\to P$、$Q'\to Q$，所求的延拓存在，当且仅当\Term{form}{形式} $d(E_\eta)$ 在 $P_0'/P_0'^0\otimes Q_0'/Q_0'^0$ 上的\Term{restriction}{限制}为零。
\SourcePage{315}
\textbf{7.3.3. \Term{symmetry}{对称性}。} 设 $E_\eta$ 是 $(P_\eta,Q_\eta)$ 的以 $\mathbf G_{m,\eta}$ 为\Term{kernel}{核}的\Term{biextension}{双扩张}，并考虑其\Term{symmetric-biextension}{对称双扩张} ${}^{s}E_\eta$（第VII讲2.7），这是 $(Q_\eta,P_\eta)$ 的以 $\mathbf G_{m,\eta}$ 为\Term{kernel}{核}的\Term{biextension}{双扩张}。则下图交换：
\[
\begin{tikzcd}[column sep=large]
\Phi\otimes\Psi\arrow[d,"\mathrm{sym}"']\arrow[dr,"d(E_\eta)"]&\\
\Psi\otimes\Phi\arrow[r,"d({}^{s}E_\eta)"']&(\mathbf Q/\mathbf Z)_k
\end{tikzcd}\tag{7.3.3.1}
\]
其中竖直\Term{arrow}{箭头}为\Term{canonical}{典范}\Term{symmetry}{对称映射}。（可将此命题与2.2.10的\Term{antisymmetry}{反对称性}命题比较。）特别地，由此可知，若 $P=Q$，且 $E_\eta$ 是 $(P_\eta,Q_\eta)$ 的一个以 $\mathbf G_{m,\eta}$ 为\Term{kernel}{核}的\Term{symmetric-biextension}{对称双扩张}，即 $E_\eta$ \Term{isomorphism}{同构}于 ${}^{s}E_\eta$，则 $\Phi\times\Phi$ 上取值于 $(\mathbf Q/\mathbf Z)_k$ 的\Term{bilinear-form}{双线性形式} $d(E_\eta)$ 是\Term{symmetric}{对称}的（并不像2.2.12那样是\Term{alternating}{交错}的！）。还应注意，若此\Term{form}{形式}为零，即 $E_\eta$ 可以延拓为 $(P,P)$ 的以 $\mathbf G_{m,S}$ 为\Term{kernel}{核}的\Term{biextension}{双扩张} $E$，则后者也必然\Term{symmetric}{对称}，因为由7.1 a），任意\Term{isomorphism}{同构} ${}^{s}E_\eta\to E_\eta$ 唯一延拓为\Term{isomorphism}{同构} ${}^{s}E\to E$。
\SourcePage{316}
\textbf{7.3.4. 将（7.1.2）型\Term{obstruction}{障碍}归结为（7.1.1）型\Term{obstruction}{障碍}。}

为明确起见，假定存在整数 $n>0$ 使 $n\Phi=0$，并设
\[
q\in Q(S)
\]
为 $Q$ 的一个\Term{section}{截面}。略用记号
\[
q_0\in\Psi(k)
\]
表示 $q$ 在 $s$ 处的值在 $\Psi(k)$ 中的像。把 $E_\eta$ 看作 $P_\eta\times Q_\eta$ 的以 $\mathbf G_{m,Q_\eta}$ 为\Term{kernel}{核}的\Term{group-extension}{扩张}，得到\Term{group-extension}{扩张}
\[
q_\eta^*(E_\eta),
\]
这是 $P_\eta$ 的以 $\mathbf G_{m,\eta}$ 为\Term{kernel}{核}的\Term{group-extension}{扩张}，从而得到\Term{obstruction}{障碍}
\[
d_q\mathrel{\mathop=^{\mathrm{dfn}}}d\bigl(q_\eta^*(E_\eta)\bigr):\Phi\longrightarrow(\mathbf Q/\mathbf Z)_k.\tag{7.3.4.1}
\]
于是，置 $d=d(E_\eta)$，便有关系
\[
d_q(x)=d(x,q_0),\tag{7.3.4.2}
\]
它对 $\Phi$ 的任意取值于 $k$-\Term{scheme}{概形} $T$ 的点 $x$ 都成立。

\textbf{7.3.4.3.} 这个公式对每个 $q\in Q(S)$ 都用\Term{pairing}{配对} $d$ 确定了 $d_q$；反过来，只要 $\Psi$ 的每个点都来自 $Q$ 的一个\Term{section}{截面}，就显然可以用各\Term{invariant}{不变量} $d_q$ 恢复这个\Term{pairing}{配对}。由 Hensel \Term{lemma}{引理}，当 $S$ \Term{strictly-local}{严格局部}时，这个条件自动成立。由于在 $k$ 的一个\Term{separable-closure}{可分闭包}上知道\Term{pairing}{配对} $d$ 就知道它本身，而由下文7.3.5.3，其形成与 $S$ 上的\Term{etale-localization}{平展局部化}交换，可见，允许换成 $S$ 的\Term{strict-localization}{严格局部化}后，确定\Term{pairing}{配对} $d$ 的问题可归结为确定各\Term{invariant}{不变量} $d_q$。
\SourcePage{317}
特别地，若 $\Psi$ 的每个点都来自 $Q$ 的一个\Term{section}{截面} $q$，就可看出，$d$ 为零（即 $E_\eta$ 延拓为 $(P,Q)$ 的以 $\mathbf G_m$ 为\Term{kernel}{核}的\Term{biextension}{双扩张} $E$），当且仅当对 $Q$ 的每个\Term{section}{截面} $q$，$d_q$ 都为零，即 $q_\eta^*(E_\eta)$ 可以延拓为 $P$ 的以 $\mathbf G_m$ 为\Term{kernel}{核}的\Term{group-extension}{扩张}。若不对 $\Psi$ 作假设，则只要让 $q$ 取遍取值于 $S$ 上\Term{etale}{平展}的\Term{local-scheme}{局部概形} $S'$ 的点，类似的判据仍成立。

\textbf{7.3.4.4.} 反过来，从7.1 a）的情形出发，取 $Q=\mathbf Z_S$ 以及由 $E$ \Term{derive}{导出}的 $(P,\mathbf Z_S)$ 的\Term{biextension}{双扩张}，就可把\Term{obstruction}{障碍}（7.1.1）解释为（7.1.2）型的\Term{obstruction}{障碍}。为说明这一点，只须在上述讨论中取 $q$ 为 $\mathbf Z_S$ 的\Term{section}{截面} $1$。

\textbf{7.3.5. \Term{base-change}{基变换} $S'\to S$ 的影响。} 假定给定\Term{trait}{离散赋值环谱}之间的一个\Term{local-morphism}{局部}\Term{dominant-morphism}{支配态射}
\[
f:S'\longrightarrow S,
\]
用 $\eta'$ 表示 $S'$ 的\Term{generic-point}{泛点}，$s'$ 表示其\Term{closed-point}{闭点}，并置 $k'=k(s')$。若 $S=\operatorname{Spec}(V)$、$S'=\operatorname{Spec}(V')$，则 $V'$ 是\Term{dominant}{支配}\Term{discrete-valuation-ring}{离散赋值环} $V$ 的\Term{discrete-valuation-ring}{离散赋值环}。置
\[
e(S'/S)=e(V'/V)=\operatorname{long}(V'/\mathfrak mV'),\tag{7.3.5.1}
\]
其中 $\mathfrak m$ 为 $V$ 的\Term{maximal-ideal}{极大理想}。现在取7.1 a）中的 $P$，令 $E_\eta$ 为 $P_\eta$ 的以 $\mathbf G_{m,\eta}$ 为\Term{kernel}{核}的\Term{group-extension}{扩张}，$P'=P\times_S S'$，并令 $E_{\eta'}'$ 为 $E_\eta$ 的\Term{inverse-image}{逆像}，它是 $P_{\eta'}'$ 的以 $\mathbf G_{m,\eta'}$ 为\Term{kernel}{核}的\Term{group-extension}{扩张}。于是有\Term{canonical-isomorphism}{典范同构}
\SourcePage{318}
\[
\Phi'\mathrel{\mathop=^{\mathrm{dfn}}}P_0'/P_0'^0\simeq
\Phi\otimes_k k'\simeq\Phi_{k'},
\]
这里采用4.1 a）的记号，因而可以考虑
\[
d(E_\eta)_{k'}:\Phi'\longrightarrow(\mathbf Q/\mathbf Z)_{k'}.
\]
于是有关系
\[
d(E_{\eta'}')=e(S'/S)\,d(E_\eta)_{k'}.\tag{7.3.5.2}
\]
同样，在7.1 b）的条件下，采用显然的记号，仍有关系（7.3.5.2）。此时 $d(E_\eta)$ 表示\Term{pairing}{配对}（7.1.2），$d(E_\eta)_{k'}$ 表示从它经\Term{base-change}{基变换} $k\to k'$ 得到的\Term{pairing}{配对}
\[
\Phi'\otimes\Psi'\longrightarrow(\mathbf Q/\mathbf Z)_{k'},
\]
而 $E_{\eta'}'$ 是通过\Term{base-change}{基变换} $\eta'\to\eta$，由 $(P_\eta,Q_\eta)$ 的以 $\mathbf G_{m,\eta}$ 为\Term{kernel}{核}的\Term{biextension}{双扩张} $E_\eta$ 得到的 $(P_{\eta'}',Q_{\eta'}')$ 的以 $\mathbf G_{m,\eta'}$ 为\Term{kernel}{核}的\Term{biextension}{双扩张}。

\textbf{7.3.5.3.} 一个有意义的情形是 $e(S'/S)=1$，即
\[
\mathfrak m'=\mathfrak mV',
\]
其中 $\mathfrak m'$ 表示 $V'$ 的\Term{maximal-ideal}{极大理想}。特别地，若 $V'$ 在 $V$ 上\Term{etale}{平展}，或者是 $V$ 的\Term{henselization}{亨泽尔化}、\Term{strict-henselization}{严格亨泽尔化}或\Term{completion}{完备化}，就是这种情形。此时，关系（7.3.5.2）简单地表明，\Term{obstruction}{障碍}（7.1.1）及其相应情形（7.1.2）的形成，与所考虑的\Term{base-change}{基变换} $S'\to S$ 交换。

\textbf{7.3.5.4.} 另一方面，\Term{generic}{一般}情形下的关系（7.3.5.2）表明：对7.1 a）中的任意\Term{group-extension}{扩张} $E_\eta$（相应地，对7.1 b）中的任意\Term{biextension}{双扩张} $E_\eta$），存在\Term{trait}{离散赋值环谱}的一个\Term{finite-morphism}{有限态射} $S'\to S$，使 $k(\eta')$ 为 $k(\eta)$ 的\Term{separable-extension}{可分扩张}，且由 $E_\eta$ 经\Term{base-change}{基变换} $\eta'\to\eta$ 得到的\Term{group-extension}{扩张}（相应地，\Term{biextension}{双扩张}）$E_{\eta'}'$ 可以在整个 $S'$ 上延拓为\Term{group-extension}{扩张}（相应地，\Term{biextension}{双扩张}）$E'$。
\SourcePage{319}
事实上，只须选择 $S'/S$，使 $e(S'/S)$ 是7.1中整数 $n>0$ 的倍数；而我们知道，总可以找到这样的\Term{extension}{扩张}，并使 $k(\eta')$ 在 $k(\eta)$ 上\Term{separable-extension}{可分}。

\begin{SGAresult}{命题}{7.4}
设 $S$ 是\Term{reduced}{约化}\Term{scheme}{概形}。

 a）设 $P$ 为 $S$ 上的\Term{smooth}{光滑}\Term{group-scheme}{群概形}，其在 $S$ 的各\Term{maximal-point}{极大点}处的\Term{fiber}{纤维}都\Term{connected}{连通}。则第VII讲1.3.4.1的自然\Term{functor}{函子}
\[
\operatorname{EXT}(P,\mathbf G_{m,S})\longrightarrow
\operatorname{TORSRIG}(P,e_P;\mathbf G_{m,S}),\tag{7.4.1}
\]
从 $P$ 的以 $\mathbf G_{m,S}$ 为\Term{kernel}{核}的\Term{group-extension}{扩张}的\Term{category}{范畴}到相对于 $P$ 的\Term{unit-section}{单位截面}作了\Term{rigidification}{刚化}的 $\mathbf G_{m,P}$-\Term{torsor}{挠子}的\Term{category}{范畴}，是\Term{fully-faithful}{充分忠实}的。若 $S$ \Term{normal}{正规}，且 $P$ 的各\Term{fiber}{纤维}都\Term{connected}{连通}，则（7.4.1）右端的一个\Term{object}{对象} $E$ 属于这个\Term{functor}{函子}的\Term{essential-image}{本质像}，当且仅当对 $S$ 的每个\Term{maximal-point}{极大点} $\eta$，$P_\eta$ 上的\Term{rigidification}{刚化}\Term{torsor}{挠子} $E_\eta$ 都来自 $P_\eta$ 的一个以 $\mathbf G_{m,\eta}$ 为\Term{kernel}{核}的\Term{group-extension}{扩张}。

 b）设 $P,Q$ 是 $S$ 上的两个\Term{smooth}{光滑}\Term{group-scheme}{群概形}，其在 $S$ 的各\Term{maximal-point}{极大点}处的\Term{fiber}{纤维}都\Term{connected}{连通}。则第VII讲2.9.1的自然\Term{functor}{函子}
\[
\operatorname{BIEXT}(P,Q;\mathbf G_{m,S})\longrightarrow
\operatorname{TORSBIRIG}(P,Q;\mathbf G_{m,S}),\tag{7.4.2}
\]
从 $(P,Q)$ 的以 $\mathbf G_{m,S}$ 为\Term{kernel}{核}的\Term{biextension}{双扩张}的\Term{category}{范畴}到相对于 $P,Q$ 的一对\Term{unit-section}{单位截面}作了\Term{birigidification}{双刚化}的 $\mathbf G_{m,P\times_S Q}$-\Term{torsor}{挠子}的\Term{category}{范畴}，是\Term{fully-faithful}{充分忠实}的。
\SourcePage{320}
若 $S$ \Term{normal}{正规}，且 $P$ 或 $Q$ 的各\Term{fiber}{纤维}都\Term{connected}{连通}，则（7.4.2）右端的一个\Term{object}{对象} $E$ 属于这个\Term{functor}{函子}的\Term{essential-image}{本质像}，当且仅当对 $S$ 的每个\Term{maximal-point}{极大点} $\eta$，$E_\eta$ 都来自 $(P_\eta,Q_\eta)$ 的一个以 $\mathbf G_{m,\eta}$ 为\Term{kernel}{核}的\Term{biextension}{双扩张}。
\end{SGAresult}
\begin{proof}
我们证明 a）；b）的证明本质相同，留给读者。对于\Term{fully-faithful}{充分忠实}性，注意，若 $X,Y$ 是\Term{reduced}{约化}\Term{scheme}{概形} $S$ 上的\Term{flat}{平坦}\Term{scheme}{概形}，且其在各\Term{maximal-point}{极大点}处的\Term{fiber}{纤维}满足4.1所列的条件，则4.1仍成立：事实上，$X\times_S Y$ 在\Term{reduced}{约化}的 $S$ 上\Term{flat}{平坦}，所以其各\Term{maximal-fiber}{极大纤维}组成的族在 $X\times_S Y$ 中\Term{schematically-dense}{概形稠密}，从而 $X\times_S Y$ 到\Term{separated}{分离} $S$-\Term{scheme}{概形} $\mathbf G_{m,S}$ 的一个\Term{morphism}{态射}，只要在各\Term{maximal-fiber}{极大纤维}上平凡，就整体平凡。于是，（7.4.1）的\Term{fully-faithful}{充分忠实}性的证明与4.2相同，且该处给出的\Term{essential-image}{本质像}的刻画也仍成立。还须证明：若 $P$ 上的\Term{rigidification}{刚化}\Term{torsor}{挠子} $E$ 的各 $E_\eta$ 来自\Term{group-extension}{扩张}，则在 $S$ \Term{normal}{正规}、$P$ 的各\Term{fiber}{纤维}都\Term{connected}{连通}的假设下，$E$ 也如此。后一个条件蕴含 $P$ 在 $S$ 上\Term{finite-presentation}{有限表示}（SGA 3，第VI$_B$讲5.5）；另一方面，\Term{fully-faithful}{充分忠实}性说明问题在 $S$ 上是\Term{local-scheme}{局部}的，故可假定 $S$ 为\Term{local-scheme}{局部}\Term{affine-scheme}{仿射概形}。通常的取\Term{limit}{极限}过程（利用 Nagata 关于\Term{reduced}{约化}\Term{finite-type}{有限型} $\mathbf Z$-\Term{algebra}{代数}的\Term{integral-closure}{整闭包}的\Term{finiteness-theorem}{有限性定理}，EGA IV，7.8）把问题归结为 $S$ 还\Term{noetherian}{诺特}的情形。此时7.1 a）表明，存在 $S$ 的一个\Term{open-subset}{开集} $U$，其补集的\Term{codimension}{余维}至少为 $2$，使得 $E_U$ 来自 $P_U$ 的一个以 $\mathbf G_{m,U}$ 为\Term{kernel}{核}的\Term{group-extension}{扩张}。换言之，所求的\Term{isomorphism}{同构}
\SourcePage{321}
\[
\pi^*(E)\simeq\operatorname{pr}_1^*(E)\operatorname{pr}_2^*(E)
\]
可以在 $(P\times_S P)_U$ 上构造。而这是\Term{normal}{正规}\Term{scheme}{概形} $P\times_S P$ 的一个\Term{open-subset}{开集}，其补集的\Term{codimension}{余维}至少为 $2$。因此，$(P\times_S P)_U$ 上所说的\Term{isomorphism}{同构}延拓到整个 $P\times_S P$，证明完毕。
\end{proof}
\begin{SGAresult}{推论}{7.5}
设 $S$ 为\Term{normal}{正规}\Term{scheme}{概形}，$P,Q$ 为 $S$ 上两个\Term{smooth}{光滑}\Term{group-scheme}{群概形}，其各\Term{maximal-fiber}{极大纤维}\Term{connected}{连通}，且 $P$ 或 $Q$ 的所有\Term{fiber}{纤维}都\Term{connected}{连通}。再假定对 $S$ 的每个\Term{maximal-point}{极大点} $\eta$，$P_\eta$ 或 $Q_\eta$ 是\Term{abelian-scheme}{阿贝尔概形}的以\Term{torus}{环面}为\Term{kernel}{核}的\Term{group-extension}{扩张}（更\Term{generic}{一般}地，假定它有一个\Term{composition-series}{合成列}，其各因子是\Term{abelian-scheme}{阿贝尔概形}、\Term{torus}{环面}或\Term{group}{群} $\mathbf G_a$）。则\Term{functor}{函子}（7.4.2）是\Term{equivalence-of-categories}{范畴等价}。
\end{SGAresult}
这确实是7.4 b）与4.7的直接推论。

\SGAsection{}{参考文献}
\begin{description}
\item[{[1]}] M. Demazure，J. Giraud，M. Raynaud，\emph{Schémas abéliens}（\Term{abelian-scheme}{阿贝尔概形}），Orsay 研讨班，1967/68，待刊。
\item[{[2]}] M. Raynaud，\emph{Faisceaux amples sur les schémas en groupes et les espaces homogènes}（\Term{group-scheme}{群概形}与\Term{homogeneous-space}{齐性空间}上的\Term{ample-sheaf}{丰沛层}），巴黎博士论文，1968（将由 Springer 在 Lecture Notes 中出版）。
\item[{[3]}] M. Rosenlicht，\emph{Some basic theorems on algebraic groups}（\Term{algebraic-group}{代数群}的若干基本\Term{theorem}{定理}），\emph{Amer. Jour. Math.}，第78卷，第2期（1956），第401--443页。
\item[{[3bis]}] M. Rosenlicht，\emph{Some rationality questions on algebraic groups}（\Term{algebraic-group}{代数群}的若干有理性问题），\emph{Annali di Matematica}，第4辑，第43号（1957），第25--50页。
\item[{[4]}] J. P. Serre，\emph{Corps locaux}（\Term{local-field}{局部域}），\emph{Act. Sci. Ind.}，1296（1962），Hermann，巴黎。
\end{description}
